% !PRINTED none
% part-title page (no running head, no printed page number): centred heading
% "Nachträge." with a horizontal rule beneath it, then subsection "I." and
% "Irreducibilität der Kreistheilungsgleichung.", followed by body text.
\begin{center}
{\Large \gesperrt{Nachträge.}}

\rule{4em}{0.4pt}

I.

\textbf{\gesperrt{Irreducibilität der Kreistheilungsgleichung.}}
\end{center}

Der \S.\,134 des ersten Bandes bedarf in seinem letzten Theile, der von der Irreducibilität der Kreistheilungsgleichung in dem allgemeinen Falle handelt, in dem der Grad der Einheitswurzel eine aus mehreren verschiedenen Primzahlen zusammengesetzte Zahl ist, einer Berichtigung, die hier nachgetragen werden soll, und wir bitten den Leser, die genannte Stelle nach dem, was hier ausgeführt wird, zu ergänzen.

Wenn $n=ab$ gesetzt wird, $a$ und $b$ relativ prim sind, und $\varrho, \alpha, \beta$ primitive $n^{\mathrm{te}}, a^{\mathrm{te}}, b^{\mathrm{te}}$ Einheitswurzeln sind, so ist $\varrho=\alpha\beta$, und es ist an der angeführten Stelle bewiesen, dass, wenn $\varrho$ Wurzel einer rationalen Gleichung $\Phi(x)=0$ ist, unter den Wurzeln dieser Gleichung $\varrho=\alpha\beta$ alle primitiven $a^{\mathrm{ten}}$ Einheitswurzeln $\alpha$ und alle primitiven $b^{\mathrm{ten}}$ Einheitswurzeln $\beta$ vorkommen müssen. Daraus folgt aber nicht, wie dort irrthümlich angenommen ist, dass darunter \gesperrt{alle} primitiven $n^{\mathrm{ten}}$ Einheitswurzeln vorkommen, weil nicht gezeigt ist, dass jedes $\alpha$ mit jedem $\beta$ combinirt vorkommt.

Zur Ausfüllung dieser Lücke möge hier zunächst ein ganz elementarer Beweis Platz finden.\footnote{\gesperrt{Arndt}, \gesperrt{Crelle}'s Journ., Bd. 56, S. 178 (1859). \gesperrt{Lebesgue}, \gesperrt{Liouville}'s Journ., 2. Ser., Bd. 4, S. 105 (1859).}

Ist $p$ eine Primzahl, so sind alle Binomialcoëfficienten für die $p^{\mathrm{te}}$ Potenz

\[
B_h^{(p)} = \dfrac{p(p-1)\ldots(p-h+1)}{1\,.\,2\ldots h},
\]

mit Ausnahme von $B_0^{(p)}$ und $B_p^{(p)}$, durch $p$ theilbare ganze Zahlen,
%CONT>>
